\documentclass[11pt]{article}

\usepackage[margin=1in]{geometry}
\usepackage{amsmath, amssymb, amsthm}
\usepackage{booktabs}
\usepackage{graphicx}
\usepackage{enumitem}
\usepackage[colorlinks=true, linkcolor=blue, citecolor=blue, urlcolor=blue]{hyperref}
\usepackage{xcolor}
\usepackage{tcolorbox}

\newtcolorbox{designbox}{colback=blue!4, colframe=blue!40, boxrule=0.5pt, arc=2pt,
  left=6pt, right=6pt, top=4pt, bottom=4pt}

\newtheorem{observation}{Observation}
\newtheorem{principle}{Principle}

\newcommand{\SI}{\mathrm{SI}}
\newcommand{\E}{\mathbb{E}}

\title{\textbf{NichePopulation: Competition-Driven Specialization and\\
Reward-Independent Capacity Assignment for Learner Populations}\\[0.4cm]
\large An Explanatory Companion: Architecture, Mechanisms, and Demonstrations}

\author{Yuhao Li\\
University of Pennsylvania\\
\texttt{li88@sas.upenn.edu}}

\date{\today}

\begin{document}

\maketitle

\begin{abstract}
Learning systems that must serve many distinct environmental regimes with
capacity-limited learners face a persistent allocation problem: which learner
handles which regime, and---when the environment is non-stationary---who
\emph{keeps} a regime that falls dormant? In practice this gap appears as three
recurring failure modes: homogenization (every learner converges to the same
behavior), coverage gaps (no learner masters rare regimes), and catastrophic
forgetting (dormant-regime knowledge is overwritten to serve active ones). To
address these failures, we present \textbf{NichePopulation}, a deliberately
minimal population-learning framework built on three tightly coupled components:
\emph{per-regime belief learning} (Thompson sampling over a shared method
inventory), \emph{niche-affinity tracking} (a canonical exponentiated-gradient
update on the regime simplex), and \emph{winner-take-all competitive exclusion}
(only the per-step winner reinforces its niche). At the core of the framework is
a single emergent property with outsized consequences: the converged
regime-to-agent assignment is \emph{structural}---independent of current task
reward---so a specialist simply idles while its niche is dormant and retains its
expertise. We show that this property, not specialization per se, is what
separates mechanisms that survive non-stationarity from mechanisms that do not:
across five capacity-allocation arms (capacity-bounded monolith, learned
Mixture-of-Experts router, fixed random niches, an EOI/CDS-style intrinsic
diversity objective, and emergent competition), every reward-chasing allocation
forgets dormant regimes and every reward-independent assignment retains them
($+71\%$ lower post-reactivation error, $p\sim10^{-35}$; robust to swapping
hard eviction for soft interference). Empirical results across six real-world
domains (145{,}294 records) and controlled synthetic tasks show that the
framework reliably self-organizes into one-per-regime specialists (mean SI
$=0.65$ from competition alone), matches purpose-built diversity and routing
baselines on spatial coverage with none of their machinery, and is structurally
immune to a failure mode of reward-driven routing. Code and data:
\url{https://github.com/HowardLiYH/NichePopulation}.
\end{abstract}

\vspace{0.3cm}
\noindent\textbf{Keywords:} emergent specialization; competitive exclusion;
multi-agent learning; catastrophic forgetting; mixture-of-experts;
reward-independent assignment

\newpage
\tableofcontents
\newpage

%==============================================================================
\section{Introduction}
%==============================================================================

\begin{quote}
\emph{``Two species with identical ecological niches cannot stably coexist.''}\\
---the competitive exclusion principle \cite{gause1934struggle,hardin1960competitive}
\end{quote}

Populations of learners that share an environment---trading strategies in a
market, sensors in a distributed network, experts inside a mixture-of-experts
model---face a coordination problem that single learners do not: how should the
population \emph{divide} a heterogeneous task space, without a central
coordinator, so that every regime of the environment is served by a competent
specialist? Despite decades of progress in multi-agent learning, three
intertwined limitations persist in standard approaches.

First, \textbf{shared objectives homogenize}. Cooperative multi-agent
reinforcement learning methods (IQL, VDN, QMIX, MAPPO)
\cite{tan1993multi,rashid2018qmix,yu2022surprising} optimize a shared task
return, and a well-documented consequence is that agents converge to
near-identical policies. In our head-to-head re-runs, all four methods reach a
Specialization Index below $0.02$ on every domain---no division of labor emerges
by default. Second, \textbf{explicit coordination is expensive}. Methods that do
produce diverse roles---diversity-MARL (ROMA, RODE, MAVEN, EOI, CDS)
\cite{wang2020roma,wang2021rode,mahajan2019maven,jiang2021eoi,li2021cds},
quality-diversity archives \cite{mouret2015illuminating}, or learned gating
networks---require auxiliary machinery: intrinsic
identifiability rewards, behavior descriptors chosen \emph{a priori}, or a
trained router. Each adds objectives, hyperparameters, and failure modes of its
own. Third, and least appreciated, \textbf{reward-driven capacity allocation
forgets}. When the environment is non-stationary and per-learner capacity is
bounded, any mechanism that allocates capacity by chasing current reward---a
single capacity-bounded model, or a Mixture-of-Experts gate trained on task
reward---systematically overwrites knowledge of regimes that are currently
dormant, and pays a full relearning cost each time a regime reactivates. The
dormant regime is precisely the one emitting no reward signal, so a reward-driven
allocator receives no signal to protect it.

To address these limitations, we introduce a cohesive design philosophy centered
on a single ecological mechanism: \textbf{competitive exclusion}. The framework,
named \textbf{NichePopulation}, is built from three integrated components:

\textbf{First}, each agent maintains \emph{per-regime method beliefs}---Beta
distributions over the success of each prediction method in each regime, acted
on by Thompson sampling. This is the agent's competence: what it knows how to do,
and where. \textbf{Second}, each agent tracks a \emph{niche affinity}
distribution over regimes, updated by the canonical exponentiated-gradient
(Hedge / multiplicative-weights) rule. This is the agent's identity: where it
belongs. \textbf{Third}, agents are coupled solely through \emph{winner-take-all
competition}: at each step, all agents act, and only the best performer
reinforces its affinity for the current regime. No communication, no shared
reward, no gate, no diversity objective. Competition makes homogeneous strategies
unstable---identical agents split the same wins, so deviating to a less-contested
regime pays---and the population partitions into one-per-regime specialists.

The key insight of this document is that the resulting assignment has a property
that purpose-built alternatives lack: \emph{it is reward-independent at
convergence}. A converged specialist holds its niche by identity, not by
moment-to-moment reward, so when its regime goes dormant the specialist simply
idles and \emph{retains} its expertise. We formalize this as a
\textbf{reward-independence principle} (Section~\ref{sec:principle}), show that
it cleanly dichotomizes five capacity-allocation mechanisms in controlled
experiments (every reward-chasing arm forgets; every reward-independent arm
retains), and demonstrate that it independently corroborates recent
mixture-of-experts theory, which finds that the gating network's updates must be
\emph{terminated} for continual-learning convergence.

In summary, this work makes the following key contributions:

\textbf{1) A reward-independence principle for retention.} We identify
reward-independence of capacity assignment as the property governing whether a
capacity-bounded system retains dormant-task knowledge under non-stationarity,
verify it five-for-five across allocation mechanisms, and explain it with an
idealized observation: dormant tasks emit no protective reward signal, so only
structural assignments can reserve capacity for them.

\textbf{2) A coordination-free specialization mechanism.} NichePopulation couples
winner-take-all competition with an exponentiated-gradient niche update and
nothing else. Competition alone (niche bonus $\lambda=0$) drives mean SI to
$0.65$ across six real domains; with the default bonus the population reaches
SI $\approx 0.99$ and one-per-regime coverage.

\textbf{3) Honest benchmarking against purpose-built baselines.} Rather than
compare only to cooperative MARL (which does not try to specialize), we compare
to an EOI/CDS-style learned-diversity objective and a learned MoE router at
matched capacity. Spatial coverage proves to be a commodity all three achieve;
competition's value there is parsimony. Retention is not a commodity: the
reward-driven router fails it decisively.

\textbf{4) A meaningfulness audit.} We show that the Specialization Index
certifies \emph{organization}, not exploitable structure (it is invariant to
regime-label shuffling and inflatable by the learning rate), and we localize
exactly when division of labor improves task accuracy---under bounded capacity
and non-stationarity---and when it does not.

%==============================================================================
\section{Related Work}
%==============================================================================

We review prior work along five dimensions closely related to NichePopulation.

\textbf{Cooperative MARL.} Independent Q-Learning, value-decomposition methods
(VDN, QMIX), and multi-agent policy gradients (MAPPO) optimize shared objectives
with centralized training. These methods are the right baselines for cooperative
control, but they do not attempt diversity, and empirically they homogenize: in
our re-runs all reach SI $\le 0.02$. We treat them as confirmatory contrast, not
as competitors on specialization.

\textbf{Diversity- and role-emergence MARL.} ROMA \cite{wang2020roma} and RODE
\cite{wang2021rode} learn role embeddings; MAVEN \cite{mahajan2019maven}
injects diverse exploration; EOI \cite{jiang2021eoi} and CDS \cite{li2021cds}
add intrinsic individuality or mutual-information objectives. These methods do produce
division of labor, at the cost of explicit diversity machinery. We benchmark
directly against an EOI/CDS-style intrinsic-reward arm and find it a strong
baseline: it matches competition on coverage for $K\ge2$ and retains dormant
regimes (its assignment is driven by an identity reward, not task reward---the
reward-independence principle again). Competition's advantage over it is
parsimony, not dominance.

\textbf{Quality-diversity optimization.} MAP-Elites \cite{mouret2015illuminating}
and novelty search \cite{lehman2011evolving} maintain archives of diverse high performers, but require behavior descriptors defined
\emph{a priori}. NichePopulation needs no descriptors: niches are discovered by
the dynamics.

\textbf{Mixture-of-experts in continual learning.} Recent theory
\cite{li2024moecl} establishes that MoE can mitigate catastrophic forgetting by
isolating tasks in experts---%
\emph{provided} experts are plentiful relative to tasks and the gating network's
updates are eventually terminated, without which continual task arrival
destabilizes routing. Our experiments occupy exactly the regime those provisos
exclude (fixed total capacity, experts forced into reuse, gate still learning)
and show that isolation then breaks down: the router inherits the monolith's
forgetting. Freezing the gate is, in our terms, making the assignment
reward-independent; competition reaches that state emergently, with no freezing
schedule to tune.

\textbf{Ecological niche theory.} The competitive exclusion principle
\cite{gause1934struggle,hardin1960competitive} and MacArthur's resource
partitioning \cite{macarthur1958population} describe how competition alone
drives species into complementary niches. NichePopulation operationalizes this
for artificial learners and inherits its most useful property: the partition,
once formed, is maintained by identity rather than by ongoing payoff, which is
what makes it robust to dormancy.

%==============================================================================
\section{The NichePopulation Framework: A Unified Architecture for
Emergent Division of Labor}
%==============================================================================

This section presents the integrated architecture. We first articulate the
design philosophy (3.1), then detail the architectural blueprint (3.2), and
finally explain the per-step dynamics that realize self-organization (3.3).

\subsection{Design Philosophy: The Competitive-Exclusion Loop}

Contemporary approaches to division of labor treat the assignment of learners to
tasks as something to be \emph{optimized}: a gate is trained, a diversity
objective is added, roles are learned. NichePopulation is founded on the opposite
principle: \textbf{the assignment should be an equilibrium, not an objective}.
We formalize this as the competitive-exclusion loop, a recursive paradigm
structured around four functions:

\begin{designbox}
\textbf{Compete}: all agents act on the current regime; one wins.\\
\textbf{Reinforce}: only the winner strengthens its affinity for that regime.\\
\textbf{Differentiate}: losers' relative affinities drift toward less-contested
regimes.\\
\textbf{Stabilize}: once each agent owns a regime, deviation is unprofitable;
the partition is a fixed point.
\end{designbox}

The key property of this loop is \emph{what it leaves behind}. Because the
converged partition is a fixed point of identity (each agent's affinity
concentrates on its niche), it no longer depends on the flow of reward. An agent
whose niche is dormant receives no wins, but it also receives no pressure to
abandon its niche---it simply idles. This is the structural source of the
retention results in Section~\ref{sec:experiments}: the loop intrinsically
unifies \emph{organizing} with \emph{remembering}.

\subsection{Architectural Blueprint: Components and Their Roles}

The loop is instantiated through three principal components per agent, plus one
population-level rule.

\textbf{1. Method-Belief Store: The Competence Layer.} Each agent $i$ maintains,
for each regime $r$ it has encountered, a Beta distribution
$\mathrm{Beta}(\beta^{+}_{i,r,m},\, \beta^{-}_{i,r,m})$ over the success rate of
each method $m$ in a shared inventory of $M$ methods. This store answers
\emph{``what works here?''} and is updated by standard Bayesian counting. In the
bounded-capacity experiments the store is explicitly finite: it holds at most
$K$ regimes, with either least-recently-used eviction (hard model) or graded
interference decay (soft model)---Section~\ref{sec:capacity}.

\textbf{2. Niche Affinity: The Identity Layer.} Each agent maintains a
distribution $\alpha_i \in \Delta^R$ over regimes---its evolving answer to
\emph{``where do I belong?''} Affinity is updated only on wins, by the canonical
exponentiated-gradient rule (Section~\ref{sec:affinity}), and read out through
the Specialization Index $\SI(\alpha) = 1 - H(\alpha)/\log R$, which is $1$ for
a one-hot specialist and $0$ for a uniform generalist.

\textbf{3. Thompson-Sampling Actor: The Decision Layer.} At each step the agent
samples $\tilde{\theta}_m \sim \mathrm{Beta}(\beta^{+},\beta^{-})$ for every
method and plays the argmax---a principled exploration--exploitation balance
requiring no tuned exploration schedule.

\textbf{4. Winner-Take-All Rule: The Population Coupling.} The \emph{only}
interaction between agents is the per-step competition: the agent with the
highest (optionally niche-shaped) reward is declared winner and is the only one
to reinforce its affinity. There is no communication channel, no shared critic,
no gate, and no diversity term. This minimality is deliberate: every emergent
property of the system must be attributable to competition, because nothing else
couples the agents.

\subsection{System Dynamics: One Step of the Loop}

The components interact in a fast per-step cycle. Given the environment in
regime $r_t$:

\begin{enumerate}[itemsep=2pt]
    \item \textbf{Act.} Every agent $i$ selects a method $m_i$ by Thompson
    sampling from its beliefs for $r_t$ (agents without competence in $r_t$ act
    uninformed) and receives raw reward $R_i$.
    \item \textbf{Score.} Each agent's competitive score adds an additive niche
    bonus centred on the uniform affinity:
    $\tilde{R}_i = R_i + \lambda\,(\alpha_{i,r_t} - 1/R)$, with $\lambda \ge 0$.
    At $\lambda = 0$ the bonus vanishes and raw reward decides.
    \item \textbf{Select winner.} $i^* = \arg\max_i \tilde{R}_i$.
    \item \textbf{Update beliefs.} Agents update their Beta counts for the
    methods they played (success/failure), within their capacity budget.
    \item \textbf{Update identity.} Only the winner applies the
    exponentiated-gradient affinity update for $r_t$ and renormalizes.
\end{enumerate}

Two timescales emerge from this single cycle without being designed in: a
\emph{fast} competence dynamic (beliefs sharpen within a regime in tens of
steps) and a \emph{slow} identity dynamic (affinities concentrate over hundreds
of steps). The slow dynamic is the one with the architectural consequence: once
identities have concentrated, the regime-to-agent map is fixed and
reward-independent, and the population behaves as a distributed, persistent
memory over the regime space.

%==============================================================================
\section{Niche Affinity and the Exponentiated-Gradient Update}
\label{sec:affinity}
%==============================================================================

\subsection{The Update Rule}

When agent $i$ wins in regime $r$ with normalized reward (quality) $q \in [0,1]$,
its affinity updates multiplicatively and renormalizes:
\begin{equation}
    \alpha_{i,r} \leftarrow \frac{\alpha_{i,r}\, e^{\eta q}}
    {\sum_{r'} \alpha_{i,r'}\, e^{\eta q\,\mathbf{1}[r'=r]}},
\end{equation}
the canonical exponentiated-gradient / Hedge / multiplicative-weights step on
the regime simplex. Three properties motivate this choice over additive
heuristics. \textbf{(i) Simplex preservation by construction:} affinities remain
a probability distribution with no clamping (an earlier additive variant of this
project required an undocumented clamp whose pathology suppressed roughly half
the specialization signal; replacing it with EG is what separates the current
V4 results from the V3 ones). \textbf{(ii) A constant effective rate:} the
multiplicative step does not slow down or speed up depending on the current
state in the way the clamped additive rule did. \textbf{(iii) A standard form:}
EG on the simplex is the textbook update for tracking the best coordinate,
making the algorithm's one learnable dynamic maximally legible. We deliberately
do \emph{not} claim Hedge's adversarial regret bound transfers to our coupled
winner-take-all setting; that would require separate analysis.

\subsection{Worked Demonstration}

Consider $R=3$ regimes, $\eta = 0.6$, and an agent at uniform affinity
$\alpha = (0.333, 0.333, 0.333)$. The agent wins three consecutive steps in
regime $1$ with quality $q = 1$:
\[
(0.333,0.333,0.333) \;\to\; (0.477,0.262,0.262) \;\to\; (0.624,0.188,0.188)
\;\to\; (0.751,0.124,0.124).
\]
Three wins move SI from $0$ to $0.39$; ten wins take it past $0.8$. The same
trajectory in reverse never occurs spontaneously: losing in regime 1 does not
push the agent away (losers do not update); it is \emph{winning elsewhere} that
re-anchors identity. This asymmetry---identity moves only on positive evidence
of fit---is what makes converged identities stable through dormancy.

\subsection{Heterogeneous Initialization and Symmetry Breaking}

Agents start near-uniform with small random perturbations. The perturbation
carries no labels and no assignment; it only breaks ties so that the first few
competitions are not decided by coin flips alone. The eventual one-per-regime
partition is an equilibrium selected by the dynamics, not an initialization
artifact: random-shuffle controls confirm the same partition statistics under
re-drawn perturbations.

%==============================================================================
\section{Competitive Exclusion and Winner-Take-All Dynamics}
%==============================================================================

\subsection{Why Strictness Matters}

Winner-take-all is not a simplification of softer competition---it is the key
structural driver of differentiation. Under proportional (softmax) credit,
identical agents can persist indefinitely: each receives a share of every
regime's reward, and hedging across niches is stable. Under winner-take-all,
identical agents split wins evenly in expectation, so each one's per-regime
reinforcement halves; an agent that deviates to an uncontested regime wins that
regime \emph{outright}. Formally (companion paper, Proposition 1), homogeneous
strategy profiles are not Nash equilibria of the winner-take-all game: a
unilateral deviation to a less-contested niche strictly increases expected
reinforcement.

\subsection{The Niche Bonus Is an Accelerator, Not the Cause}

The additive shaping term $\lambda(\alpha_{i,r_t} - 1/R)$ advantages an agent
competing in a regime it already prefers. Setting $\lambda = 0$ removes it
entirely---and specialization still emerges, with mean SI $0.65$ across the six
real domains (every domain $> 0.49$, versus $0.13$ for a random baseline). The
$\lambda$-sweep is flat and high ($\SI > 0.98$) across $\lambda \in [0.2,0.4]$
and declines at $0.5$, where over-locked winners stop contesting secondary
regimes and leave some niches under-filled. The mechanism therefore needs no
shaping to work; shaping only sharpens it.

\subsection{What the Population Looks Like After Convergence}

On every real domain, the converged population is a clean partition: each agent
holds a primary niche with affinity $\ge 0.95$, every regime is owned by at
least one agent, and the population collectively uses $86\%$ of the shared
method inventory (specialists adopt the methods suited to their niche). The
assignments are interpretable by inspection---a property no deep MARL baseline
in our comparison offers.

%==============================================================================
\section{Bounded Capacity and Memory Models}
\label{sec:capacity}
%==============================================================================

Specialization only matters when no single learner can do everything. We model
this with an explicit per-agent capacity $K$: each agent can hold beliefs for at
most $K$ of the $R$ regimes. Two memory models implement the constraint.

\subsection{Hard Model: LRU Slots}

The belief store holds $K$ regime entries; learning an uncached regime evicts
the least-recently-used entry. Eviction is total: the evicted regime's beliefs
reset to the uninformed prior, and re-encountering it costs a full relearning
period. This is the cleanest possible model of ``capacity binds.''

\subsection{Soft Model: Interference Decay}

A natural objection is that real learners do not evict discretely---knowledge
\emph{interferes}. The soft model removes eviction entirely: every learning
update on the active regime decays the pseudo-counts of the agent's other held
regimes by a factor that grows with how far the agent's learned footprint
exceeds $K$. Dormant regimes are never refreshed, so they decay smoothly toward
the prior---the interference account of catastrophic forgetting, with graded
rather than all-or-nothing loss. Every retention result in
Section~\ref{sec:experiments} is reproduced under both models; the headline
separation is a property of \emph{who gets updated}, not of how forgetting is
discretized.

\subsection{The Five Capacity-Allocation Arms}

All experiments compare five mechanisms at matched per-agent capacity $K$:

\begin{enumerate}[itemsep=2pt]
    \item \textbf{Monolith}: one agent with capacity $K$ serves everything.
    Allocation chases the active regime by necessity.
    \item \textbf{MoE router}: $N$ capacity-$K$ experts behind a learned gate
    ($\epsilon$-greedy routing; gate trained on routed-expert quality).
    Allocation chases current reward by design.
    \item \textbf{Random diversity}: $N$ agents with fixed, random capacity-$K$
    niches. Allocation is frozen---fully reward-independent, but coverage is
    accidental.
    \item \textbf{Learned diversity (EOI/CDS-style)}: all agents receive an
    intrinsic identifiability reward $\log p(i \mid r)$ on top of task quality;
    the identified owner learns the regime. Allocation follows an
    \emph{identity} signal, not task reward.
    \item \textbf{NichePopulation (ours)}: winner-take-all competition;
    allocation is an emergent, converged identity.
\end{enumerate}

This five-arm design is deliberate: the arms differ in \emph{what signal drives
capacity assignment}, which is exactly the axis the next section's principle
predicts should matter.

%==============================================================================
\section{The Reward-Independence Principle}
\label{sec:principle}
%==============================================================================

\subsection{Why a Reward-Driven Router Forgets}

Consider a regime $r$ that goes dormant for an extended interval while an expert
$e$ holding $r$ receives $D$ learning updates from other, active regimes.

\begin{observation}[Dormant regimes receive no protective signal]
Under the hard capacity model, $e$ retains $r$ only if fewer than $K$ distinct
regimes are used during the interval; under the soft model, $r$'s retained
strength decays as $(1-\rho)^D \to 0$. Crucially, a dormant regime emits no
steps and hence no reward gradient: the gate update
$\Delta\,\mathrm{gate}[r][e] \propto (q_e - \tfrac12)$ is identically zero
throughout dormancy, so the router receives no signal that would lead it to
reserve $e$ for $r$. Reactivation therefore incurs the full relearning cost with
probability $\to 1$ as dormancy grows.
\end{observation}

The asymmetry is the crux: \emph{routing and capacity protection draw on the
same signal, and the regime most needing protection is the one producing none.}
A converged competitive (or identity-driven) assignment escapes the trap because
it is reward-independent at steady state: the specialist for $r$ simply idles
through dormancy, applying zero foreign updates to its own store ($D = 0$ for
its niche), and retains it structurally---no signal required.

\begin{principle}[Reward-independent assignment]
\label{prin:ria}
Under bounded per-agent capacity and non-stationarity, a capacity-allocation
mechanism retains dormant-task knowledge if and only if its regime-to-learner
assignment is independent of current task reward.
\end{principle}

\subsection{The Five-Arm Dissociation}

Principle~\ref{prin:ria} makes a sharp prediction across the five arms of
Section~\ref{sec:capacity}, verified in full by the experiments:

\begin{center}
\small
\begin{tabular}{llcc}
\toprule
\textbf{Arm} & \textbf{Assignment signal} & \textbf{Reward-indep.?} &
\textbf{Retains?} \\
\midrule
Monolith & current active regime & no & no \\
MoE router & current task reward & no & no \\
Random diversity & none (frozen) & yes & yes \\
Learned diversity & intrinsic identity reward & yes & yes \\
NichePopulation & converged identity & yes & yes \\
\bottomrule
\end{tabular}
\end{center}

Five for five, retention tracks reward-independence---not specialization, not
population size, not learning sophistication. Among the retaining arms,
competition is the most parsimonious: random diversity retains but cannot
guarantee coverage at $K = 1$ (unowned regimes are possible); learned diversity
retains but requires the auxiliary intrinsic-reward machinery; gate-freezing
(the prescription of recent MoE continual-learning theory \cite{li2024moecl})
retains but requires deciding \emph{when} to freeze. Competition needs none of these---the
reward-independent assignment is the equilibrium the dynamics converge to.

\subsection{Relation to MoE Continual-Learning Theory}

Recent theory on mixture-of-experts in continual learning \cite{li2024moecl}
proves that the gating network's updates must be \emph{terminated} after sufficient training for the
system to converge, and that MoE's forgetting-mitigation holds when experts are
plentiful or can be added. Read through Principle~\ref{prin:ria}, both findings
are the same statement: a stabilized (frozen) gate is a reward-independent
assignment, and plentiful experts remove the forced reuse that makes
reward-chasing destructive. Our experiments cover the complementary regime the
theory excludes---fixed total capacity, forced reuse ($K \ll R$), gate still
learning---and find exactly the failure the principle predicts. The two lines of
work corroborate each other from opposite directions: regression theory
prescribes freezing; population dynamics exhibits a mechanism that arrives at
the frozen state emergently.

%==============================================================================
\section{Experiments and Evaluation}
\label{sec:experiments}
%==============================================================================

\subsection{Experimental Settings}

\textbf{Domains.} Six real-world domains with verified public data
(145{,}294 records): cryptocurrency (Bybit; 8{,}766 daily bars), commodities
(FRED; 5{,}630), weather (Open-Meteo; 9{,}105), solar irradiance (Open-Meteo;
116{,}834 hourly), urban traffic (NYC TLC; 2{,}879 hourly), and air quality
(Open-Meteo; 2{,}880 hourly). Each domain defines 4--6 regimes and 5 tailored
prediction methods. Controlled synthetic regime-champion streams (each regime
has a distinct optimal method, with a tunable separation parameter) isolate
mechanisms that real data confounds.

\textbf{Baselines.} Homogeneous (best single method), random selection,
cooperative MARL (IQL, VDN, QMIX, MAPPO) at matched agent counts and training
budgets, and---for the capacity experiments---the five-arm comparison of
Section~\ref{sec:capacity}.

\textbf{Metrics.} Specialization Index (organization), task prediction error
(performance), post-reactivation error (retention, measured in a 15-step probe
window after each regime reactivation). Protocol: 20--30 seeded trials per
condition, $95\%$ confidence intervals, one-sided Welch tests with
Holm--Bonferroni robustness checks for sweep families.

\subsection{Results on Emergent Self-Organization}

Under the default configuration, every domain converges to near-perfect
organization (mean SI $= 0.992$; all domains $> 0.98$), versus $\le 0.02$ for
all cooperative MARL baselines---a qualitative, not quantitative, gap. The
critical ablation sets $\lambda = 0$: competition alone yields mean SI $= 0.65$
(every domain $> 0.49$), $5.1\times$ the random baseline, confirming that the
niche bonus accelerates but does not cause specialization. Populations use
$86\%$ of the method inventory on average; division of labor extends to
\emph{how} agents predict, not just \emph{where}.

We treat these SI results as a diagnostic, not a headline. A battery of
controls shows SI is invariant to regime-label shuffling and inflatable by the
learning rate: high SI certifies organization, not exploitable structure. The
performance-relevant results follow.

\subsection{Results on Coverage Under Bounded Capacity}

With per-agent capacity bounded ($K < R$), the specialized population beats the
capacity-matched monolith decisively (synthetic $+73\%$ at $K{=}1$; real traffic
$+22\%$), and beats the equal-capacity \emph{random}-diversity population
specifically at tight capacity (synthetic $+57\%$, $p<10^{-16}$; traffic
$+8.7\%$, $p<10^{-3}$ at $K{=}1$), where random assignment leaves coverage gaps
that competitive exclusion provably fills. Against the purpose-built arms the
honest finding is parity: the learned-diversity objective matches competition
for $K \ge 2$ (and exceeds it on some real domains), and the MoE router matches
it on the synthetic task at every $K$. A method-overlap sweep maps the boundary:
competition's edge over learned diversity grows monotonically with method
exclusivity, from $-4.8\%$ (overlapping methods) to $+29.3\%$ (fully exclusive).
\emph{Spatial coverage is a commodity}; competition's value there is parsimony
and sample-efficiency on noisy real data (it beats the router on traffic at
$K{=}1$ by $+5.1\%$), not dominance.

\subsection{Results on Retention Under Non-Stationarity}

Regimes cycle (a sliding window of $W{=}3$ of $R{=}6$ regimes is active per
epoch), so every regime goes dormant and reactivates. This is where the five
arms dissociate exactly as Principle~\ref{prin:ria} predicts:

\begin{center}
\small
\begin{tabular}{lcc}
\toprule
\textbf{Arm ($K{=}1$, hard model)} & \textbf{Post-reactivation error} &
\textbf{Retains?} \\
\midrule
Monolith & $1.015$ & no \\
MoE router & $0.928$ & no \\
Random diversity & $0.603$ & partially (coverage gaps) \\
Learned diversity & $0.322$ & yes \\
NichePopulation & $\mathbf{0.283}$ & yes \\
\bottomrule
\end{tabular}
\end{center}

The population beats the capacity-matched monolith by $+33\%$ overall and
$+71\%$ in the post-reactivation window ($p<10^{-36}$ at $K{=}3$), and the
router---which matched competition on spatial coverage---fails here decisively
($0.928$ vs.\ $0.283$ at $K{=}1$, $p\sim10^{-35}$; the gap stays significant
through $K{=}4$, $p<10^{-12}$), closing only as $K \to R$. The specialized
population is essentially flat in $K$: even $K{=}1$ agents collectively form a
complete persistent memory.

\subsection{Ablation: Robustness to the Memory Model}

Re-running the entire sweep under the soft interference model (no eviction at
all) reproduces the dissociation: the monolith still forgets (post-reactivation
error $0.85$--$0.92$ at $K \le 3$), the specialized population still retains
($\approx 0.25$, a $+70\%$ gap, $p\sim10^{-40}$), and the router remains
significantly worse than the population ($+36\%$ at $K{=}1$, $p\sim10^{-8}$;
$+16\%$ at $K{=}3$). As the Observation predicts, the router degrades more
gracefully under graded interference than under hard eviction---but never closes
the gap until $K \to R$. The separation is a property of reward-driven
allocation, not of any particular forgetting mechanism.

\subsection{Ablation: The Honest Audit}

Three controls scope all of the above. \textbf{(i)} SI is invariant to
regime-label shuffling: it measures organization, not structure. \textbf{(ii)}
An oracle structure-diagnostic shows four of six domains have a globally
dominant method (no exploitable regime structure); on the two structured domains
(traffic, solar), causal leakage-free regime labels beat shuffled labels by
$+15.4\%$ and $+12.8\%$. \textbf{(iii)} With \emph{unlimited} capacity, a single
regime-conditioned monolith matches or beats the population: division of labor
is a task-accuracy lever \emph{specifically} under bounded capacity and
non-stationarity, not in general. We consider this audit as much a contribution
as the positive results: it states exactly where the mechanism's value lives.

%==============================================================================
\section{Conclusion and Future Work}
%==============================================================================

We presented NichePopulation, a minimal framework in which winner-take-all
competition over a niche-affinity simplex makes a population of capacity-bounded
learners self-organize into one-per-regime specialists---with no communication,
shared rewards, gates, or diversity objectives. The framework's deepest property
is that its converged assignment is \emph{reward-independent}: specialists hold
niches by identity, idle through dormancy, and thereby retain what reward-driven
allocators---including the standard learned MoE router---systematically forget.
Across five capacity-allocation mechanisms, retention tracks
reward-independence of assignment five-for-five; across six real-world domains
and controlled synthetic tasks, competition matches purpose-built diversity and
routing baselines on spatial coverage with none of their machinery, and
separates from them decisively on retention ($+71\%$ post-reactivation,
$p\sim10^{-35}$, robust to the memory model). An honest audit bounds the claim:
with slack capacity and stationary regimes, a regime-conditioned monolith
suffices; division of labor pays specifically when capacity binds and the world
shifts.

Several directions remain promising. First, a function-approximation
demonstration: the reward-independence principle predicts the router's failure
is architecture-agnostic, and validating it with gradient-trained experts on a
standard continual-learning benchmark would convert the prediction into
evidence. Second, hybrid allocators: a router augmented with an explicit
memory-preservation or load-balancing term could in principle reserve idle
experts for dormant regimes---our results specify \emph{what objective such a
router would need}, and testing one would complete the design space. Third,
scaling the population: decentralized competition over hundreds of regimes,
where the one-per-regime equilibrium must coexist with agent churn and partial
observability. Finally, deployment in real elastic systems---trading-strategy
pools, sensor networks, and expert pools inside large models---where the
parsimony of competition (no gate to train, no schedule to freeze, no diversity
coefficient to tune) is most valuable.

\bibliographystyle{plain}
\bibliography{references}

\end{document}
